\section{Panorama bibliométrico e corpus da revisão}
\label{sec:panorama}

O núcleo final reúne 104 registros publicados entre 2011 e 2026. O período de 2019 a 2026 concentra 88 registros, correspondentes a 84,6\% do núcleo, com maior frequência em 2025, quando foram identificados 25 estudos. O Gráfico~\ref{fig:temporal} apresenta a distribuição anual.

\begin{figure}[H]
\centering
\includegraphics[width=0.95\linewidth]{figuras/figura09_distribuicao_temporal_nucleo_final_104.png}
\captiongrafico{Distribuição temporal do núcleo final}
\label{fig:temporal}
\end{figure}

A deduplicação preservou todas as bases em que cada estudo foi recuperado. Entre os 104 registros únicos, 52 foram encontrados apenas na Scopus, 41 na Scopus e na Web of Science, cinco apenas na Web of Science, três nas três bases, dois na Scopus e no Crossref e um apenas no Crossref. Por essa razão, a contagem por base representa proveniência e admite respostas múltiplas: 98 registros possuem origem Scopus, 49 possuem origem Web of Science e seis possuem origem Crossref.

A tipologia documental foi harmonizada para evitar que rótulos equivalentes fossem apresentados como categorias distintas. Os rótulos \textit{Journal}, \textit{JOUR} e \textit{journal-article} foram reunidos como artigo de periódico. \textit{Conference Proceeding} e \textit{CPAPER} foram reunidos como trabalho em evento. \textit{Book Series} e \textit{Book} foram reunidos como livro ou série de livro. A Tabela~\ref{tab:basetipo} apresenta as duas dimensões separadamente.

\begin{table}[H]
\centering
\caption{Distribuição do núcleo final por base de origem e tipo documental}
\label{tab:basetipo}
\footnotesize
\begin{tabularx}{\textwidth}{>{\RaggedRight\arraybackslash}p{3.7cm} >{\raggedleft\arraybackslash}p{1.8cm} >{\raggedleft\arraybackslash}p{1.8cm} Y}
\toprule
Categoria & Registros & \% dos 104 & Observação \\
\midrule
\multicolumn{4}{l}{\textbf{A. Base de origem, com respostas múltiplas}} \\
Scopus & 98 & 94,2 & Base bibliográfica principal. \\
Web of Science & 49 & 47,1 & Base bibliográfica independente. \\
Crossref & 6 & 5,8 & Fonte complementar de metadados. \\
\addlinespace
\multicolumn{4}{l}{\textbf{B. Tipo documental harmonizado, categorias exclusivas}} \\
Artigo de periódico & 79 & 76,0 & \textit{Journal}, \textit{JOUR} e \textit{journal-article}. \\
Trabalho em evento & 15 & 14,4 & \textit{Conference Proceeding} e \textit{CPAPER}. \\
Livro ou série de livro & 10 & 9,6 & \textit{Book Series} e \textit{Book}. \\
\bottomrule
\end{tabularx}
\fonteautor
\end{table}

